\chapter{Literature and Theoretical Background}
\label{ch:literature}

This chapter develops the conceptual framework that motivates the empirical analysis. \Cref{sec:asset-prices-fundamental-value} introduces the present-value model of asset pricing and defines a speculative bubble formally, alongside the main interpretations of rapid asset-price increases that coexist in the literature. \Cref{sec:technology-speculation} reviews the historical pattern linking technological innovation to speculative excess, and notes that the formal test used in this thesis has previously detected canonical historical episodes. \Cref{sec:productivity-paradox} turns to the productivity paradox literature, which shows why the economic benefits of transformative technologies can lag market valuations by years or decades, and why this lag is relevant for AI. \Cref{sec:detecting-bubbles} presents the PSY econometric procedure used to detect explosive price dynamics, together with its main applications. \Cref{sec:bubbles-gap-contribution} draws these strands together and sets out the gap this thesis addresses.

\section{Asset prices and fundamental value}
\label{sec:asset-prices-fundamental-value}

The starting point for any analysis of speculative excess is the present-value model of asset pricing. Under standard assumptions, the price of an asset is determined by

\begin{equation}
P_t = \sum_{i=0}^{\infty} \left( \frac{1}{1+r_f} \right)^i E_t \left( D_{t+i} + U_{t+i} \right) + B_t,
\label{eq:present-value-bubble}
\end{equation}

\noindent where $P_t$ is the asset price, $r_f$ is the risk-free rate, $D_t$ denotes the payoff (such as dividends), $U_t$ represents unobservable factors, and $B_t$ is a bubble component. When $B_t = 0$, the asset is priced by fundamentals alone in this representation. The empirical implication used by the PSY literature is that the relevant price or valuation series should not exhibit persistent explosive behaviour under the no-bubble null. A speculative bubble arises when $B_0 > 0$, in which case the bubble component satisfies the submartingale property

\begin{equation}
E_t(B_{t+1}) = (1+r_f)\, B_t.
\label{eq:bubble-submartingale}
\end{equation}

\noindent When $B_t$ is present, it dominates price dynamics and transmits explosive behaviour to the observed price series \citep{phillipsShiYu2015a}.

The decomposition in \cref{eq:present-value-bubble} formalises what a bubble is, but it does not settle when one should be invoked. Three views about how to interpret rapid asset-price increases coexist in the literature, and they imply different prior beliefs about whether explosive dynamics observed during a technology boom should be attributed to $B_t$ at all.

The strongest rejection of the bubble interpretation comes from the efficient market hypothesis \citep{fama1970}. In its strict form, asset prices at all times fully reflect available information, so sustained mispricings cannot occur: any deviation from fundamental value would be arbitraged away. Under this view, the divergence in returns between AI-exposed firms and the rest of the index, if it exists, reflects a rational updating of expectations about future earnings capacity and should not be interpreted as evidence of mispricing without additional evidence.

A more nuanced rational benchmark is provided by \citet{pastorVeronesi2009}. They show that the arrival of a major new technology raises uncertainty about future productivity, which in turn raises both the level and the volatility of asset prices for firms exposed to the new technology. As information about the technology's true economic impact is gradually revealed, prices adjust, sometimes sharply downward. The mechanism produces price paths that are observationally similar to bubble episodes --- sharp run-ups followed by corrections --- but it generates them without any departure from rationality. Whether a specific episode reflects rational learning of this kind or a genuine bubble component cannot be determined from price data alone; the two hypotheses generate similar reduced-form patterns and must be distinguished, if at all, with additional structure or auxiliary evidence.

A third view, due to \citet{shiller2000}, holds that financial markets are subject to waves of irrational exuberance, in which investor psychology, media-driven amplification, and feedback loops drive prices above any level justified by careful analysis of fundamentals. Under this view, the bubble component $B_t$ is not only theoretically admissible but empirically relevant, and explosive dynamics in asset prices during technology booms are natural candidates for speculative excess unless auxiliary evidence points to the rational-learning alternative.

These three views set the interpretive boundaries for the empirical work that follows. The PSY procedure introduced in \cref{sec:detecting-bubbles} provides a formal test for explosive price dynamics, but it does not by itself discriminate between speculative excess and rational repricing under uncertainty. A finding of explosive dynamics is consistent with both \citet{shiller2000} and \citet{pastorVeronesi2009}; a finding of no explosive dynamics is consistent with \citet{fama1970} and with the rational-learning view in episodes where the underlying uncertainty has already been resolved. The cross-sectional design developed in this thesis (\cref{sec:bubbles-gap-contribution}) is intended to sharpen, not resolve, this identification challenge. If explosive dynamics are stronger in portfolios most closely tied to the AI narrative, and remain visible relative to broad-market and non-AI benchmarks, the evidence is more consistent with an AI-specific speculative component than with a purely aggregate market boom.

\section{Technology and speculation: a recurring pattern}
\label{sec:technology-speculation}

The link between technological innovation and speculative excess is one of the most persistent features of financial history. The financial-crises literature documents a recurring sequence: a displacement, typically a new technology or financial innovation, triggers a period of rising asset prices, which attracts additional capital, loosens credit conditions, and produces a self-reinforcing cycle of enthusiasm that eventually overshoots any plausible estimate of fundamental value \citep{kindlebergerAliber2015}. The cycle resolves, in the canonical account, through a sharp correction in which prices fall back toward (or below) fundamentals as the original premise that motivated the rally is reassessed.

The specific technologies change, but the pattern is remarkably stable. In the 1920s, the Radio Corporation of America (RCA) became a focal point for investor enthusiasm about radio broadcasting. Its stock rose from roughly \$94 to \$505 between early 1928 and September 1929, substantially outpacing the company's actual sales performance, before collapsing in the crash of October 1929 \citep{brunerMiller2019,sirkin1975}. The dot-com boom of the late 1990s followed a similar arc. Cisco Systems, the largest networking-equipment company on the Nasdaq, peaked at \$132 in March 2000 before declining to \$16 within a year \citep{hirschey2001}. In both cases the underlying technology proved economically transformative over a longer horizon; the speculative dynamic was not that the technology failed but that its near-term valuation had run far ahead of its near-term cash-flow capacity.

These historical episodes are not only narrative parallels but also detectable in formal tests. Applying the GSADF procedure to the S\&P~500 price-dividend ratio, \citet{phillipsShiYu2015a} date-stamp explosive episodes that include the 1928--1929 run-up and the 1995--2001 dot-com bubble. This precedent motivates the use of the same class of recursive right-tailed tests in the present thesis, but the empirical object is different: rather than testing only a broad aggregate, the analysis below asks whether explosive dynamics vary systematically with AI exposure and market concentration.

Two features of the historical record matter for the empirical strategy that follows. First, the price dynamics are not gradual: rapid acceleration is followed by a sharp correction. This is the signature pattern that recursive right-tailed unit-root tests of the kind developed by \citet{phillipsWuYu2011} and \citet{phillipsShiYu2015a} are designed to detect. Second, the episodes share an underlying structural feature with the current AI boom: a new technology with genuinely large but highly uncertain long-run economic potential, in which the market must price an asset whose fundamental value depends on outcomes that have not yet been realised. Whether the response to that uncertainty is rational learning, as in \citet{pastorVeronesi2009}, or speculative excess, as in \citet{shiller2000}, the resulting price paths can be similar in their reduced-form properties; this is the identification challenge raised in \cref{sec:asset-prices-fundamental-value} and revisited in \cref{sec:bubbles-gap-contribution}.

\section{The productivity paradox and the lag problem}
\label{sec:productivity-paradox}

A central difficulty in evaluating whether AI valuations are justified is that transformative technologies often take far longer to generate measurable economic returns than initial market enthusiasm implies. This problem was formalised by \citet{brynjolfsson1993} as the ``productivity paradox'' of information technology: despite massive investment in computing power from the 1970s onward, aggregate productivity statistics showed little improvement. As Robert Solow famously remarked, computers could be seen everywhere except in the productivity statistics \citep{brynjolfsson1993}.

\citet{brynjolfsson1993} identifies several possible explanations for this gap: measurement errors in output statistics, lags between investment and the organisational changes required to exploit new technology, redistribution of rents rather than creation of new value, and mismanagement. The most analytically important of these is the lag hypothesis. Drawing on an analogy to the electrification of factories, \citet{david1990} shows that major productivity gains from electric power did not materialise for roughly two decades, because the full benefits required redesigning production processes around the new technology.

This pattern appears to be repeating with AI. Drawing on coordinated surveys of approximately 6{,}000 firms across the United States, the United Kingdom, Germany, and Australia conducted between November 2025 and January 2026, \citet{yotzovEtAl2026} document that around 70\% of firms now use AI in some form, with adoption skewed toward larger and more productive firms, but that the realised productivity impact over the preceding three years has been modest, with 89\% of firms reporting no impact and an average across all firms of approximately 0.3\%. Executives expect AI to raise productivity by approximately 1.4\% over the next three years, a figure that is positive but modest relative to the scale of corporate investment and equity-market repricing documented in \cref{ch:introduction}. The gap between expected and realised productivity impact --- substantial expectations, modest current contribution --- is the contemporary version of the pattern \citet{brynjolfsson1993} and \citet{david1990} describe for earlier technologies.

This gap matters directly for the present analysis. The firms in the High AI Exposure group are, by construction, the firms that were already discussing AI as a meaningful part of their business before the post-2022 boom. If the productivity gains that justify these firms' current valuations follow the lag pattern documented for previous transformative technologies and now apparently for AI itself, the gap between current investment and future payoff is precisely the space in which the bubble component $B_t$ can develop.

\section{Detecting bubbles: from theory to econometrics}
\label{sec:detecting-bubbles}

The conceptual framework developed in \cref{sec:asset-prices-fundamental-value,sec:technology-speculation,sec:productivity-paradox} establishes that speculative excess is plausible in the current AI boom but does not, on its own, determine whether it is actually present. Answering that question requires an econometric procedure that can test whether prices exhibit the explosive dynamics characteristic of a non-zero $B_t$. This section introduces the PSY procedure, which provides a formal test for such dynamics, while leaving the economic source of the explosiveness to be interpreted with auxiliary evidence and cross-sectional structure.

The methodology was developed in a sequence of papers by Phillips and co-authors. \citet{phillipsWuYu2011} (PWY) proposed an initial recursive right-tailed augmented Dickey--Fuller (ADF) test, the SADF, applied to the 1990s Nasdaq and date-stamping bubble origination in July 1995. \citet{phillipsShiYu2015a,phillipsShiYu2015b} (PSY) then generalised the test to allow detection of \emph{multiple} bubbles within a single sample, an extension that matters when the sample period contains more than one candidate episode. The PSY framework is built as follows. Let $y_t$ be a dividend-adjusted price series. Under the null hypothesis of no bubble, the series satisfies

\begin{equation}
y_t = dT^{-\eta} + y_{t-1} + \varepsilon_t, \qquad \varepsilon_t \sim \text{iid}(0, \sigma^2),
\label{eq:psy-null}
\end{equation}

\noindent where $d$ is a constant, $T$ is the sample size, and $\eta$ is a localising coefficient that accommodates an asymptotically negligible drift when $\eta > 1/2$. Under this specification, the series follows a unit-root process rather than a mildly explosive process. This is the statistical no-bubble benchmark used by the test; it is not a complete structural model of fundamentals.

The alternative hypothesis of bubble presence is given by mildly explosive dynamics \citep{phillipsMagdalinos2007}:

\begin{equation}
y_t = \theta_T \, y_{t-1} + \varepsilon_t,
\label{eq:psy-alternative}
\end{equation}

\noindent where the autoregressive parameter $\theta_T$ exceeds unity:

\begin{equation}
\theta_T = 1 + \frac{c}{T^{\alpha}}, \qquad \alpha \in (0,1), \quad c > 0.
\label{eq:mildly-explosive-root}
\end{equation}

\noindent When $\theta_T > 1$, the price series exhibits explosive behaviour, growing at a rate that exceeds what a unit-root process would imply. The PSY procedure tests for this explosive behaviour using a sequence of right-tailed ADF statistics computed over recursive subsamples. The backward sup ADF (BSADF) statistic for each observation is

\begin{equation}
\text{BSADF}_t(\tau_0) = \sup_{\tau_1 \in [0,\, t - \tau_0]} \left\{ \text{ADF}_{\tau_1}^{t} \right\},
\label{eq:bsadf}
\end{equation}

\noindent where $\tau_0$ is the minimum window size and $\text{ADF}_{\tau_1}^{t}$ is the ADF $t$-statistic estimated on the subsample from $\tau_1$ to $t$. The generalised sup ADF (GSADF) statistic is then

\begin{equation}
\text{GSADF}(\tau_0) = \sup_{\tau \in [\tau_0,\, 1]} \left\{ \text{BSADF}_\tau(\tau_0) \right\}.
\label{eq:gsadf}
\end{equation}

Under the null hypothesis, the GSADF statistic has a known limiting distribution. Rejection of the null in favour of the explosive alternative provides evidence of bubble-like dynamics in the price series. The procedure also allows for date-stamping: a bubble is identified as originating at the first observation where the BSADF statistic exceeds the corresponding critical-value sequence, and as terminating when it falls back below.

The methodology has been applied widely. \citet{phillipsShi2018} use it to study financial-bubble implosion in equity markets, \citet{etienneIrwinGarcia2015} apply it to commodity futures, \citet{caspiKatzkeGupta2018} to crude oil prices, and \citet{pavlidisEtAl2016} to housing markets across multiple metropolitan areas. The Federal Reserve Bank of Dallas uses the procedure to monitor global housing-market exuberance in real time \citep{phillipsShi2020}, and software implementations are available in EViews and R.

The most direct methodological precedent for the present thesis is \citet{baselePhillipsShi2025}, who apply the PSY procedure to the Nasdaq index and to the seven individual stocks of the Magnificent Seven over January 2010 to January 2025. They use the price-to-index ratio as the test variable to control for market fundamentals and the wild bootstrap of \citet{harveyEtAl2016} for critical values that account for heteroskedasticity. Their findings indicate persistent speculative dynamics in the Nasdaq index across both halves of the sample, and explosive behaviour in six of the seven Mag-7 stocks during the December 2022 to January 2025 subperiod. The empirical strategy developed in \cref{ch:empirics} adopts the same procedure and the same bootstrap framework, but applies them to value- and equal-weighted portfolios of S\&P~500 firms grouped by their pre-2023 AI exposure, together with relative ratios and concentration baskets. The conceptual motivation for that extension is set out in \cref{sec:bubbles-gap-contribution}.


\section{Gap and contribution}
\label{sec:bubbles-gap-contribution}

The literatures reviewed in \cref{sec:asset-prices-fundamental-value,sec:technology-speculation,sec:productivity-paradox,sec:detecting-bubbles} form four distinct strands. The first establishes a formal definition of a speculative bubble through the present-value model, in which a non-zero component $B_t$ satisfying the submartingale condition in \cref{eq:bubble-submartingale} drives explosive dynamics in the observed price series \citep{phillipsShiYu2015a}. The second documents that technological innovation has repeatedly been associated with asset-price dynamics that appear to run ahead of realised fundamentals, while also providing a rational benchmark in which genuine uncertainty about a transformative technology can generate bubble-like paths ex post \citep{kindlebergerAliber2015,pastorVeronesi2009}. The third shows that the economic benefits of transformative technologies typically lag their market impact by years or decades, creating a structural gap between investor expectations and realised returns in which $B_t$ can develop \citep{brynjolfsson1993,david1990}. The fourth provides a formal statistical test for explosive dynamics, together with a date-stamping procedure that identifies the origination and termination of candidate episodes \citep{phillipsWuYu2011,phillipsShiYu2015a,phillipsShiYu2015b}.

The intersection of these strands has so far been examined mainly at the aggregate and mega-cap level. \citet{baselePhillipsShi2025} apply the PSY procedure to the Nasdaq and to the seven individual stocks of the Magnificent Seven over January 2010 to January 2025, finding evidence of speculative bubbles in the index and in six of the seven stocks during the December 2022 to January 2025 subperiod. Their analysis establishes that bubble-like dynamics are detectable in the headline aggregates of the current AI boom, but it does not address three questions that are central for this thesis: whether the dynamics extend beyond a small set of mega-cap firms, whether they vary systematically with a disclosure-based measure of AI engagement, and whether they remain visible after broad-market dynamics are netted out.

The contribution of this thesis is therefore to extend the PSY framework along a cross-sectional and relative-pricing dimension. Rather than testing only a single index or a small set of pre-identified mega-cap stocks, the analysis groups S\&P~500 firms by their pre-2023 AI exposure score (\cref{ch:data}) and applies the PSY procedure to value- and equal-weighted portfolios, broad-market ratios, within-market High AI / No AI ratios, and concentration baskets. This design matters because absolute index rejections alone are not sufficient to establish an AI-specific interpretation: a broad bull market, changes in discount rates, risk appetite, or general technology-sector repricing can also generate explosive dynamics. The sharper test is whether explosiveness is stronger in High-AI portfolios than in Moderate- or No-AI portfolios, whether it survives relative to the S\&P~500 and to the No-AI group, and whether it becomes stronger when the portfolio is narrowed to the firms most closely associated with the AI narrative.

The empirical sections that follow should therefore be read as evidence on the cross-sectional structure of bubble-like dynamics, not as a structural decomposition of prices into fundamentals and mispricing. A pattern in which High-AI ratios and AI-leader concentration baskets reject more strongly than broad-market and placebo ratios would support the interpretation that the current episode contains an AI-specific speculative component. A pattern in which all groups reject similarly would instead point toward market-wide forces. Intermediate patterns are informative because they show how much of the apparent AI boom is broad and how much is concentrated among firms most exposed to, or most identified with, the AI narrative. The empirical strategy that operationalises this design is developed in \cref{ch:empirics}.
